% !TEX TS-program = pdflatex
% David - Tekninen harjoittelija CV (suomi)

\documentclass[a4paper,11pt]{article}

\IfFileExists{cv-style.sty}{\usepackage{cv-style}}{\usepackage{../../shared/cv-style}}

\cvlang{fi}
\cvrole{Tekninen harjoittelija}

\begin{document}
\small
\setlength{\cvlabelw}{3.2cm}
\titlespacing{\section}{0pt}{0.55em}{2pt}

\begin{cvheader}
\par{\centering
		{\Huge \textsc{David Adegbola}}\\[0.4em]
		{\Large Tekninen harjoittelija}
	\bigskip\par}

\hrule
\vspace{0.5em}
\cvline[2.4cm]{\cvt{location}:}{Joensuu, Suomi}
\cvline[2.4cm]{\cvt{phone}:}{+358 46 800 7210}
\cvline[2.4cm]{\cvt{email}:}{\href{mailto:adavidadegbola@gmail.com}{adavidadegbola@gmail.com}}
\cvline[2.4cm]{Linkit:}{\faGithub{} \href{https://github.com/david-adegbola}{github.com/david-adegbola}}
\end{cvheader}

\section{\cvt{summary}}
\begin{tabular}{p{0.9\textwidth}}
 Hei! Olen David, 24-vuotias päätoiminen tutkinto-opiskelija Suomessa: opiskelen kestävää metsäbiotaloutta Itä-Suomen yliopistossa. Opiskelin tietotekniikkaa VAMKissa ja suoritin OND-tutkinnon tietojenkäsittelystä. Ohjelmoin Pythonilla ja C:llä perustasolla ja teen pieniä sulautettuja projekteja.\\\\

 Rakennan myös työkaluja tekoälyagenteilla ja kirjoitan niille tarkistukset. UPM Kaukaan sellutehtaalla työnantaja arvioi työni kiitettäväksi ja kuvasi minua halukkaaksi oppimaan ja ottamaan vastuuta. Puhun sujuvaa englantia. Suomen kielen taitoni on vielä perustasolla, ja opiskelen sitä aktiivisesti.
\end{tabular}

\section{\cvt{education}}
\cvline{2026--nykyhetki:}{B.Sc, kestävä metsäbiotalous, Itä-Suomen yliopisto (UEF), päätoiminen opiskelija}
\cvline{2024--2026:}{B.Eng-opinnot, tietotekniikka, VAMK; kesken}
\cvline{2019--2022:}{OND, tietojenkäsittely, The Polytechnic, Ibadan, Nigeria}

\section{\cvt{skills}}
\cvline{Ohjelmointi:}{Python ja C (perustaso). Raspberry Pi 3 -lämpötilaloggeri Linuxilla (kurssiprojekti, tehty tekoälyagentin avulla); Arduinon liikennevalo-ohjain C/C++:lla, simuloitu (tein itse)}
\cvline{GitHub ja tekoäly:}{Repot, GitHub Actions ja Pages; Claude Code; säännöt ja tarkistukset tekoälyagenteille}
\cvline{Tietotekniikka:}{Microsoft Office, myös Excel; Microsoft Active Directory (VAMKin opinnot)}
\cvline{Haluan oppia:}{SQL, Power BI, GitLab, testiautomaatio, sulautettujen järjestelmien testaus}
\cvline{\cvt{cards}:}{Työturvallisuuskortti}
\cvline{\cvt{languages}:}{Englanti (sujuva); suomi (perustaso, opiskelen aktiivisesti); yoruba (äidinkieli)}

\section{\cvt{experience}}
\begin{cventry}{tammi 2026--nykyhetki}{\textbf{Tekoälykehittäjä (omat projektit)}}
  Rakennan työkaluja tekoälyagenttien avulla (Claude Code): kirjoitan määrittelyn ja agentin säännöt (faktat, kielletyt väitteet, tarkistukset) ja tarkistan tuloksen.\par
  Rakensin oman CV-järjestelmän, joka julkaisee CV:t GitHub Pagesissa.\par
  Rakensin JobDeskin: metsäalan urapalvelu (Python, SQLite), joka vertaa ilmoituksia CV:hen.
  \sotagrow{ai-automation,claude-code,agent-rules,github-actions,automation}
\end{cventry}
\begin{cventry}{kesä 2026--elo 2026}{\textbf{Founder in Residence}, Vaasa Startup Factory, Vaasa: kehitimme tiimin kanssa hetkin.app-palvelun.}
\end{cventry}
\begin{cventry}{elo 2025--loka 2025}{\job{Asentaja}{Axioma Servisas UAB}}
  UPM Kaukaan sellutehdas, Lappeenranta: sähkösuodattimen uudelleenrakennushanke.\par
  Noudatin tiukkoja turvallisuusohjeita, pidin työmaan siistinä ja seurasin työkaluja ja materiaaleja.\par
  Työnantajan arvio: kiitettävä. Vastuuntuntoinen ja tunnollinen työntekijä, halukas oppimaan ja ottamaan vastuuta.
  \sotagrow{työturvallisuus,teollisuustyömaa,tiimityö,dokumentointi,luotettavuus}
\end{cventry}
\begin{cventry}{marras 2024--elo 2026}{\textbf{AV-vapaaehtoinen}, Vaasa City Church, Vaasa: hoidin ääni- ja videotekniikan joka viikonloppu.}
\end{cventry}
\begin{cventry}{huhti 2025--elo 2025}{\textbf{Lehdenjakaja}, Finnbox ja SSM, Vaasa: jakelureitit yksin aikaisin aamulla.}
\end{cventry}
\begin{cventry}{syys 2016--syys 2020}{Myynti- ja asiakaspalvelutyötä Nigeriassa (Desire Feeds, Bunmi Pastries).}
\end{cventry}

\section{\cvt{references}}
\cvt{refsavailable}

\end{document}
